% Chapter 6 — GPU Architecture, CUDA Programming, and Maximizing Occupancy  (v2: narrative edition)
% 53 pages; Part I follows the BOOK section order (SM internals -> threads/warps/blocks/grids -> limits -> PTX);
% 4 presenter intuition pages at 7, 9, 14, 15. Page numbers differ from v1. Use chapter06-lecture-script-v2-zh.md.
% v2 changes: book-plan opening (p2: SIMT -> CUDA -> memory -> roofline), question-driven roadmap (p3), worked-example p14, takeaway titles,
% unified workshop metaphor, occupancy caveat at p16, part-transition bridges, Plain-English takeaway bars.
% Compile: pdflatex chapter06-presentation-v2.tex  (run twice)
\documentclass[aspectratio=169]{beamer}

% ====== Presenter info: EDIT ME ======
\newcommand{\PresenterName}{Tianqin Cai}          % <-- 改成你的名字
\newcommand{\PresenterTitle}{Applied Scientist} % <-- 改成你的头衔
\newcommand{\PresenterOrg}{AWS}
% =====================================

\usetheme{Madrid}
\definecolor{navy}{RGB}{18,47,90}
\definecolor{kwblue}{RGB}{0,56,132}
\definecolor{okgreen}{RGB}{80,130,20}
\usecolortheme[named=navy]{structure}
\setbeamercolor{frametitle}{bg=navy,fg=white}
\setbeamercolor{title}{bg=navy,fg=white}
\setbeamercolor{block title}{bg=navy,fg=white}
\setbeamercolor{block body}{bg=gray!12,fg=black}
\setbeamercolor{block title example}{bg=okgreen,fg=white}
\setbeamercolor{block body example}{bg=okgreen!10,fg=black}
\setbeamertemplate{navigation symbols}{}
\setbeamertemplate{blocks}[rounded][shadow=true]

\usepackage{booktabs}
\usepackage{listings}
\lstset{
  language=C++,
  basicstyle=\ttfamily\fontsize{6.5}{7.6}\selectfont,
  keywordstyle=\color{kwblue}\bfseries,
  commentstyle=\color{gray}\itshape,
  stringstyle=\color{okgreen},
  showstringspaces=false,
  breaklines=true,
  columns=fullflexible
}

\newcommand{\kw}[1]{\textbf{\textcolor{kwblue}{#1}}}
\newcommand{\pe}[1]{\par{\setlength{\fboxsep}{2.5pt}\noindent\colorbox{navy!9}{\parbox{\dimexpr\textwidth-5pt\relax}{\footnotesize\itshape Plain English: #1}}}}
\newcommand{\bridge}[1]{{\footnotesize\itshape\textcolor{gray}{#1}}\par}
\newcommand{\refmark}[1]{\textsuperscript{[#1]}}

\title[AI Systems Performance Engineering]{AI Systems Performance Engineering}
\subtitle{Chapter 6 --- GPU Architecture, CUDA Programming, and Maximizing Occupancy\\{\small Part II structure preview (draft)}}
\author[\PresenterName\ \ \PresenterTitle\ \ \PresenterOrg]{\PresenterName\\\PresenterTitle\\\PresenterOrg}
\institute[]{\footnotesize Internal Training for Applied Scientists on AI Infrastructure}
\date{August 15, 2026}

\begin{document}

% ================================================================ 1
\begin{frame}
  \titlepage
\end{frame}

% ================================================================ 2
\begin{frame}[fragile]{Part II --- Anatomy of a CUDA Kernel: Device Side}
  \bridge{The hardware wants many ready warps --- now: how does CUDA code create them?}
  \begin{columns}[t]
    \column{0.52\textwidth}
\begin{lstlisting}[emph={blockIdx,blockDim,threadIdx,idx},emphstyle=\color{kwblue}\bfseries]
__global__ void myKernel(float* input, int N) {
    // unique global index for this thread
    int idx = blockIdx.x * blockDim.x + threadIdx.x;
    if (idx < N) {          // bounds check!
        input[idx] *= 2.0f; // in-place: double it
    }
}
\end{lstlisting}
    \begin{itemize}\small
      \item \kw{\texttt{\_\_global\_\_}}: runs on the device, callable from the host.
      \item Built-ins: \texttt{blockIdx} (which block am I), \texttt{blockDim} (block size),
            \texttt{threadIdx} (my rank in the block).
    \end{itemize}
    \column{0.44\textwidth}
    \begin{itemize}
      \item The kernel is written from the perspective of \kw{one thread};
            \texttt{idx} carves out its slice of the data.
      \item CUDA compiles it into device code executed by \kw{thousands or millions}
            of lightweight threads in parallel.
      \item Launch syntax (host): \kw{\texttt{<<<blocksPerGrid, threadsPerBlock>>>}} ---
            the two core parameters of \emph{every} kernel invocation.
    \end{itemize}
  \end{columns}
  \pe{write the job description for \emph{one} worker; CUDA prints a million copies, each with a different row number.}
\end{frame}

% ================================================================ 3
\begin{frame}[fragile]{Host Side: The Six-Step Data Flow}
  \begin{columns}[t]
    \column{0.55\textwidth}
\begin{lstlisting}
const int N = 1'000'000;
float *h_input = nullptr, *d_input = nullptr;   // h_=host, d_=device
cudaMallocHost(&h_input, N * sizeof(float));    // 1) pinned host alloc
for (int i = 0; i < N; ++i) h_input[i] = 1.0f;  //    init
cudaMalloc(&d_input, N * sizeof(float));        //    device alloc
cudaMemcpy(d_input, h_input, N * sizeof(float),
           cudaMemcpyHostToDevice);             // 2) H2D copy
const int threadsPerBlock = 256;
const int blocksPerGrid =                       //    = 3,907 for 1M
    (N + threadsPerBlock - 1) / threadsPerBlock;
myKernel<<<blocksPerGrid, threadsPerBlock>>>(d_input, N); // 3) launch
cudaDeviceSynchronize();                        // 4) wait for device
cudaMemcpy(h_input, d_input, N * sizeof(float),
           cudaMemcpyDeviceToHost);             // 5) D2H copy
cudaFree(d_input); cudaFreeHost(h_input);       // 6) cleanup
\end{lstlisting}
    \column{0.42\textwidth}
    \begin{itemize}
      \item Naming convention: \kw{\texttt{h\_}} = host pointer, \kw{\texttt{d\_}} = device pointer
            --- used throughout the book.
      \item \kw{\texttt{cudaMallocHost}} = \kw{pinned} (page-locked) memory ---
            prerequisite for async copies to truly overlap later.
      \item Book note: \emph{not fully optimized yet} --- this is the simple, complete
            \kw{template} the rest of the book improves on.
    \end{itemize}
  \end{columns}
\end{frame}

% ================================================================ 4
\begin{frame}{Why Pass N? The Single-Thread Perspective}
  \begin{block}{The book's answer}
    ``A CUDA kernel function is designed to work \kw{inside of a single thread}, alongside thousands
    of other threads, \kw{on a partition of the input data}. As such, N defines the size of the
    partition that this particular kernel will process.''
  \end{block}
  \begin{itemize}
    \item A CPU function could just inspect the array size. A kernel \kw{cannot} ---
          it sees a raw pointer and must be told where its world ends.
    \item Combined with \texttt{blockIdx}/\texttt{blockDim}/\texttt{threadIdx}, N lets every
          element be processed \kw{cleanly and uniquely}, in parallel, across many SMs.
    \item This is the core mental shift from CPU to GPU programming: you don't write
          ``process the array,'' you write ``process \kw{my} element --- if it exists.''
  \end{itemize}
  \pe{each of the million copied workers gets the same memo; N is the line that says where the job ends.}
\end{frame}

% ================================================================ 5
\begin{frame}{The Bounds Check --- and Lazily Surfacing Errors}
  \begin{columns}[t]
    \column{0.5\textwidth}
    \begin{block}{Worked example: N = 63}
      \small Scheduler assigns \kw{2 warps} (64 threads). Warp 1 handles elements 0--31 --- fine.
      Warp 2 expects 32--63 \kw{plus one extra thread}; without \texttt{if (idx < N)} it reads
      past the array $\Rightarrow$ \kw{\texttt{cudaErrorIllegalAddress}}.
    \end{block}
    \begin{itemize}\small
      \item Book rule: bounds checks are everywhere in CUDA; ``if you don't see it,
            \kw{understand why it's not there}.''
    \end{itemize}
    \column{0.47\textwidth}
    \begin{itemize}
      \item Kernels run \kw{asynchronously, with no per-thread exceptions}: an illegal access sets a
            \kw{global fault flag} for the whole launch.
      \item The host only sees it at the \kw{next sync or CUDA API call} --- errors surface
            \kw{lazily}.
      \item Practice: \kw{\texttt{cudaGetLastError()} + \texttt{cudaDeviceSynchronize()}}
            right after launches, so faults are caught where they happen.
    \end{itemize}
  \end{columns}
  \pe{the GPU doesn't call you when something breaks --- it leaves a note you only find the next time you check the mailbox.}
\end{frame}

% ================================================================ 6
\begin{frame}{Configuring Launch Parameters: The Recipe in Code}
  \begin{block}{The recipe (why 256? --- justified in Part I's worked example)}
    \kw{\texttt{threadsPerBlock = 256}} (= 8 warps);\quad
    \kw{\texttt{blocksPerGrid = (N + 255) / 256}} --- integer division rounds \kw{up},
    so \kw{every element is covered}.
  \end{block}
  \begin{itemize}
    \item The rounding trick generalizes: \kw{\texttt{(N + B - 1) / B}} for any block size \texttt{B}.
    \item The \kw{last block may be partially full} --- exactly why the
          \kw{bounds check} (p.~5) exists.
    \item Book tip for Blackwell: consider \kw{256--512} threads per block.
    \item Treat 256 as the \kw{starting point}, not the answer:
          \kw{profile} 128 / 512 variants before committing.
  \end{itemize}
  \pe{256 is the ``medium coffee'' of block sizes --- almost never the wrong first order; refine later if profiling says so.}
\end{frame}
% ================================================================ 7
\begin{frame}[fragile]{2D and 3D Kernel Inputs}
  \begin{columns}[t]
    \column{0.52\textwidth}
\begin{lstlisting}
__global__ void my2DKernel(float* input,
                           int width, int height) {
    int x = blockIdx.x * blockDim.x + threadIdx.x;
    int y = blockIdx.y * blockDim.y + threadIdx.y;
    if (x < width && y < height) {  // 2D bounds check
        int idx = y * width + x;    // flatten to 1D
        input[idx] *= 2.0f;
    }
}
// host: dim3 threads(16,16);  // 256 threads, 2D
//       dim3 blocks((W+15)/16, (H+15)/16);
\end{lstlisting}
    \column{0.44\textwidth}
    \begin{itemize}
      \item Natural fit for \kw{images}: e.g., a $1{,}024 \times 1{,}024$ matrix processed by
            $16 \times 16$ blocks (still 256 threads).
      \item Same pattern generalizes to 3D with \kw{\texttt{dim3(x, y, z)}} ---
            volumetric data maps directly onto the thread hierarchy.
      \item Book note: most of the book uses \kw{1D or 2D (tiled)} launches;
            in 1D, plain ints beat \texttt{dim3}.
    \end{itemize}
  \end{columns}
  \pe{same recipe, two axes: each worker gets a (row, column) badge instead of a single row number.}
\end{frame}

% ================================================================ 8
\begin{frame}[fragile]{Allocate Asynchronously: Streams and Memory Pools}
  \begin{columns}[t]
    \column{0.52\textwidth}
\begin{lstlisting}
cudaStream_t s;
cudaStreamCreateWithFlags(&s, cudaStreamNonBlocking);

float* d_buf = nullptr;
cudaMallocAsync(&d_buf, N * sizeof(float), s);
myKernel<<<blocks, threads, 0, s>>>(d_buf, N);
cudaFreeAsync(d_buf, s);  // deferred until s finishes
\end{lstlisting}
    \begin{itemize}\small
      \item \kw{Free waits only for stream \texttt{s}} --- no global
            \texttt{cudaDeviceSynchronize}, no cross-stream sync.
    \end{itemize}
    \column{0.45\textwidth}
    \begin{itemize}\small
      \item \texttt{cudaMalloc}/\texttt{cudaFree} are \kw{synchronous and expensive}:
            full-device sync + OS calls (\texttt{mmap}/\texttt{ioctl}), kernel-space context switches.
      \item \kw{\texttt{cudaMallocAsync}/\texttt{cudaFreeAsync}} draw from a per-device
            \kw{memory pool}: recycled buffers, no OS round trips, \kw{less fragmentation}
            over long training loops.
      \item Non-blocking streams avoid legacy \kw{default-stream barriers} (book tip; multistream
            overlap in Chapter 11).
    \end{itemize}
  \end{columns}
  \pe{keep a truck fleet in the lot and grab keys --- don't buy and return a truck for every delivery.}
\end{frame}

% ================================================================ 9
\begin{frame}{Tuning the Pool; PyTorch's Caching Allocator}
  \begin{columns}[t]
    \column{0.48\textwidth}
    \begin{block}{Pool knobs}
      \begin{itemize}\small
        \item \kw{\texttt{cudaMemPoolAttrReleaseThreshold}} (via
              \texttt{cudaMemPoolSetAttribute}): how much reserved memory the pool keeps
              before releasing to the OS.
        \item \kw{\texttt{cudaMemPoolTrimTo}}: proactively return memory.
        \item Trade-off: total \kw{footprint} vs.\ \kw{fragmentation}.
      \end{itemize}
    \end{block}
    \column{0.48\textwidth}
    \begin{block}{The PyTorch connection}
      \begin{itemize}\small
        \item PyTorch's \kw{caching allocator} (\texttt{PYTORCH\_ALLOC\_CONF}, formerly
              \texttt{PYTORCH\_CUDA\_ALLOC\_CONF}) is the same idea: reuse GPU memory,
              avoid a synchronous \texttt{cudaMalloc} per new tensor per iteration.
      \end{itemize}
    \end{block}
  \end{columns}
  \medskip
  \begin{itemize}
    \item Verdict from the book: one-off buffers --- blocking \texttt{cudaMalloc} is fine;
          \kw{allocation-heavy loops --- async + pools} give more consistent performance
          and higher throughput.
  \end{itemize}
\end{frame}

\end{document}
